\section{رمزنگاری و امضاهای دیجیتال}

\subsection{مبانی رمزنگاری در شبکه اقتضایی خودرویی}

پیام‌های ایمنی در شبکهٔ اقتضایی خودرویی\LTRfootnote{\lr{Vehicular Ad Hoc Network (VANET)}} اطلاعات حساسی ندارند؛ بنابراین محرمانگی برای آن‌ها الزامی نیست و تأکید اصلی بر اصالت و یکپارچگی پیام است؛ رایا و هوکسو صریحاً بیان می‌کنند که پیام ایمنی به احراز اصالت نیاز دارد اما به رمزنگاری نیاز ندارد \cite{SecuringVehicularAd}. در همین حال، مرور فارسیمدان و همکاران می‌نویسد که اطلاعات تبادل‌شده در ارتباطات خودرو با همه‌چیز\LTRfootnote{\lr{Vehicle-to-Everything (V2X)}} به‌طور کلی عمومی است و محرمانگی ممکن است بالاترین اولویت نباشد، جز در مورد اطلاعات مربوط به حریم خصوصی سرنشینان \cite{farsimadanReviewSecurityChallenges2025}. پیام‌های این شبکه‌ها کوچک و اغلب رد و بدل می‌شوند و شبکه سیار با سرعت بالا و توپولوژی پویا دارد \cite{farsimadanReviewSecurityChallenges2025}؛ از این رو الزام دسترس‌پذیری نیازمند استفاده از سازوکارهای رمزنگاری سبک و کم‌هزینه است \cite{farsimadanReviewSecurityChallenges2025}.

در معماری پیشنهادی رایا و هوکسو، قالب پیام شامل پیام اصلی، امضای دیجیتال آن با کلید خصوصی فرستنده و گواهی صادرشده از مرجع گواهی‌دهی\LTRfootnote{\lr{Certificate Authority (CA)}} است؛ برچسب زمانی برای تازگی پیام به کار می‌رود و عدد تصادفی (\lr{nonce}) که نیاز به دست‌دهی دارد، نامطلوب شمرده می‌شود \cite{SecuringVehicularAd}. الزامات امنیتی این شبکه‌ها شامل احراز اصالت، سازگاری داده، دسترس‌پذیری، انکارناپذیری، حریم خصوصی و محدودیت‌های زمانی سخت به‌خاطر سرعت بالای خودروها است \cite{SecuringVehicularAd} و ارتباطات با اولویت بالا نباید تأخیری بیش از چند ده میلی‌ثانیه تجربه کنند \cite{farsimadanReviewSecurityChallenges2025}. نکتهٔ مهم دیگر آن است که سازوکارهای رمزنگاری عمدتاً برای دفاع در برابر تهدیدهای خارجی طراحی می‌شوند و در برابر اقدامات امنیتی پیشگیرانهٔ داخلی مانند جعل موقعیت و تولید هشدار غلط توسط خودروهای مجاز آسیب‌پذیرند \cite{farsimadanReviewSecurityChallenges2025}.

\subsection{انواع امضاهای دیجیتال}

بر پایهٔ این مبانی، انواع مختلف امضای دیجیتال در منابع مرورشده مورد بحث قرار می‌گیرند. آنچه پنج منبع این گزارش پشتیبانی می‌کنند، امضای مبتنی بر هویت و امضای گروهی (در پروتکل \lr{GSIS}) و امضای مبتنی بر منحنی بیضوی است؛ دو نوع دیگر که در ادامه آمده، در هیچ‌یک از این منابع بررسی نشده‌اند.

\subsubsection{امضای مبتنی بر هویت}

نخستین نوعی که منابع به آن پرداخته‌اند، امضای مبتنی بر هویت\LTRfootnote{\lr{Identity (ID)-Based Signature}} است. پروتکل \lr{GSIS}\LTRfootnote{\lr{GSIS}} بر پایهٔ ترکیب امضای گروهی و امضای مبتنی بر هویت بنا شده و نیازمندی‌های طراحی خاص امنیت و حفظ حریم خصوصی را برای ارتباط میان دستگاه‌های ارتباطی مختلف در شبکهٔ اقتضایی خودرویی شناسی می‌کند \cite{linGSISSecurePrivacyPreserving2007}. هدف این پروتکل تأمین امنیت و حریم خصوصی همراه با ردیابی‌پذیری هر خودرو است، به‌طوری‌که ارسال‌کننده در حالت عادی ناشناس می‌ماند، اما مرجع مجاز در صورت بروز اختلاف می‌تواند هویت واقعی او را آشکار کند \cite{linGSISSecurePrivacyPreserving2007}. کارایی، اثربخشی و قابلیت اعمال این پروتکل در سناریوهای کاربردی مختلف و سیستم‌های جاده‌ای گوناگون با شبیه‌سازی گسترده بررسی شده است \cite{linGSISSecurePrivacyPreserving2007}. در مرور فارسیمدان و همکاران نیز ترکیبی از شبه‌نام و امضای مبتنی بر هویت در میان راهکارهای هویت درج شده است \cite{farsimadanReviewSecurityChallenges2025}.

نیمهٔ دوم این ترکیب، امضای گروهی\LTRfootnote{\lr{Group Signature}} است. در مرور فارسیمدان و همکاران، طرح‌های امضای گروهی به‌عنوان یکی از دفاع‌ها در برابر دستکاری پیام (همراه با همبستگی داده و چالش و پاسخ) و در برابر تحلیل ترافیک معرفی شده‌اند \cite{farsimadanReviewSecurityChallenges2025}؛ همچنین تصدیق سبک در ارتباط خودرو با زیرساخت با کلید گروهی صادرشده از مرجع قابل اعتماد در این مرور درج شده است \cite{farsimadanReviewSecurityChallenges2025}.

\subsubsection{امضای تجمیعی}

نوع بعدی، امضای تجمیعی\LTRfootnote{\lr{Aggregate Signature}}، است که هیچ‌یک از پنج منبع مرورشده در این گزارش به آن نمی‌پردازد؛ بنابراین در اینجا تعریف، طرح یا عددی برای آن ارائه نمی‌شود. تنها سازوکار مرتبطی که در این منابع برای کاهش سربار تأیید امضا بررسی شده، تأیید دسته‌ای امضاها است: تصدیق همکارانهٔ دسته‌ای امضاها به‌عنوان یکی از دفاع‌ها در برابر حملهٔ مرد میانی\LTRfootnote{\lr{Man-in-the-Middle (MiM)}} در مرور فارسیمدان و همکاران درج شده است \cite{farsimadanReviewSecurityChallenges2025}، و رایا و هوکسو پیشنهاد می‌کنند که پیام فقط در صورت مرتبط بودن تأیید شود و کلیدی که پیش‌تر تأیید شده دوباره تأیید نشود \cite{SecuringVehicularAd}.

\subsubsection{امضای بدون گواهی}

همچنین امضای بدون گواهی\LTRfootnote{\lr{Certificateless Signature}} در هیچ‌یک از منابع پنج‌گانهٔ این گزارش بررسی نشده است و هیچ تعریف یا طرحی در این بخش ارائه نمی‌شود. رویکردهای مدیریت کلیدی که این منابع واقعاً بررسی می‌کنند، دو گروه‌اند: زیرساخت کلید عمومی مبتنی بر گواهی دیجیتال \cite{SecuringVehicularAd} و امضای مبتنی بر هویت \cite{linGSISSecurePrivacyPreserving2007}.

\subsubsection{امضای مبتنی بر منحنی بیضوی}

در مقابل دو نوع نامذکور، امضای مبتنی بر منحنی بیضوی\LTRfootnote{\lr{Elliptic Curve Cryptography (ECC)}} تنوعی است که منابع با جزئیات عددی پشتیبانی می‌کنند. رایا و هوکسو انتخاب این رمزنگاری را تحلیل کرده‌اند: کلیدی به طول ۲۲۴ بیت (۲۸ بایت) و امضایی به طول ۵۶ بایت، تقریباً معادل \lr{RSA}\LTRfootnote{\lr{RSA}} با کلید ۲۰۴۸ بیت، و سربار رمزنگاری حدود ۱۴۰ بایت \cite{SecuringVehicularAd}. در مقایسهٔ تحلیلی آن‌ها، که تجربی نیست و صرفاً نشان‌دهنده است، کلید جفتی حتی با چند خودرو سربار بیشتری از رمزنگاری مبتنی بر منحنی بیضوی دارد و کلید گروهی حدود ۵۴٪ از سربار پیام را کم می‌کند، اما تعداد پیام بیشتر و پیچیدگی پروتکل گروهی را می‌طلبد \cite{SecuringVehicularAd}. نتیجهٔ این تحلیل آن است که امضای دیجیتال با وجود اینکه کارایی آن جای بهبود دارد، مناسب‌ترین و قابل‌اتکاترین راه‌حل است \cite{SecuringVehicularAd}. در مرور فارسیمدان و همکاران نیز ترکیب رمزنگاری مبتنی بر منحنی بیضوی با رمزنگاری پایلیه\LTRfootnote{\lr{Paillier}} در یکی از طرح‌های حریم خصوصی مرورشده درج شده است \cite{farsimadanReviewSecurityChallenges2025}.

\subsection{زیرساخت کلید عمومی}

پس از مرور انواع امضا، زیرساختی که صدور و مدیریت گواهی و کلید را بر عهده دارد نقشی اساسی در این سازوکارها دارد. مدیریت اعتماد خودسازمان‌یافته (مانند \lr{PGP}\LTRfootnote{\lr{PGP}}) به‌دلیل مسائل مسئولیت کافی نیست و کلیدها باید توسط مرجع مورد اعتماد صادر شوند؛ بنابراین زیرساخت کلید عمومی\LTRfootnote{\lr{Public Key Infrastructure (PKI)}} لازم است \cite{SecuringVehicularAd}. در این معماری که به زیرساخت کلید عمومی خودرویی\LTRfootnote{\lr{Vehicular Public Key Infrastructure (VPKI)}} معروف است، هر پیام امضای دیجیتال و گواهی \lr{CA} دارد \cite{SecuringVehicularAd}. در مرور فارسیمدان و همکاران، دفاع پیشنهادی در برابر ربودن نشست\LTRfootnote{\lr{Session Hijacking}} بر پایهٔ همین زیرساخت و مرجع قابل اعتماد است: واحد کنارجاده‌ای\LTRfootnote{\lr{Roadside Unit (RSU)}} هویت خودرو را نزد مرجع بررسی و سپس کلید را با آن به اشتراک می‌گذارد \cite{farsimadanReviewSecurityChallenges2025}.

در این معماری، دستگاه مقاوم در برابر دست‌کاری\LTRfootnote{\lr{Tamper-Proof Device (TPD)}} اسرار را نگه می‌دارد و پیام‌های خروجی را امضا می‌کند، دارای باتری و ساعت مستقل است که هنگام عبور از ایستگاه پایهٔ مورد اعتماد هم‌گام می‌شود و حسگرهای دست‌کاری دارد که در صورت دست‌کاری همهٔ کلیدها را پاک می‌کنند \cite{SecuringVehicularAd}. این دستگاه برای شرایط جاده بیش از حد حساس است و گران است (نمونهٔ \lr{IBM 4758} هزاران دلار هزینه دارد)، در حالی که ماژول سکوی مورد اعتماد\LTRfootnote{\lr{Trusted Platform Module (TPM)}} با قیمت چند ده دلار در برابر حملهٔ نرم‌افزاری مقاوم است اما در برابر دست‌کاری سخت‌افزاری نه \cite{SecuringVehicularAd}. پیام‌های مرتبط با مسئولیت همراه امضا و گواهی برای بررسی‌های بعدی در ثبت‌کنندهٔ دادهٔ رویداد\LTRfootnote{\lr{Event Data Recorder (EDR)}} ذخیره می‌شوند که به جعبه‌سیاه هواپیما شبیه است \cite{SecuringVehicularAd}.

مدیریت کلید شامل هویت الکترونیکی یکتا و قابل تأیید رمزنگارانه است که دولت (پلاک الکترونیکی\LTRfootnote{\lr{Electronic License Plate (ELP)}}) یا سازنده (شمارهٔ شاسی الکترونیکی\LTRfootnote{\lr{Electronic Chassis Number (ECN)}}) صادر می‌کند \cite{SecuringVehicularAd}. کلیدها هنگام ثبت خودرو بارگذاری می‌شوند، کلیدهای ناشناس به‌طور دوره‌ای (مثلاً در معاینهٔ سالانه) تجدید می‌شوند و کلید عمومی \lr{CA} از پیش بارگذاری می‌شود؛ \lr{CA} یا مراجع حمل‌ونقل دولتی است که در مرز مناطق بازگواهی‌دهی می‌کنند تا زنجیرهٔ گواهی لازم نباشد، یا سازندگان خودرو \cite{SecuringVehicularAd}.

جفت‌کلید ناشناس\LTRfootnote{\lr{Anonymous Key Pair}} جفت کلیدی است که \lr{CA} آن را تأیید می‌کند اما هیچ اطلاعاتی دربارهٔ هویت واقعی خودرو (یعنی \lr{ELP} آن) ندارد و با آن رابطهٔ عمومی ندارد؛ با این حال این ناشناسی برای مقاصد مسئولیت مشروط است و خودرو مجموعه‌ای از این کلیدها دارد تا ردیابی نشود \cite{SecuringVehicularAd}. تمام شناسه‌های خودرو، به‌ویژه آدرس‌های \lr{MAC}\LTRfootnote{\lr{MAC}} و \lr{IP}\LTRfootnote{\lr{IP}}، باید در طول زمان تغییر کنند، چون با ثبت پیام‌های یک کلید می‌توان فرستنده را تا کشف هویتش ردیابی کرد \cite{SecuringVehicularAd}. توان افشای هویت باید میان چند مرجع توزیع شود؛ به‌عنوان مثال، پلیس بدون اجازهٔ قاضی نتواند هویت را بازیابی کند. برای دسترسی به پایگاه‌دادهٔ میان \lr{ELP} و کلیدهای ناشناس از تسهیم راز شامیر\LTRfootnote{\lr{Shamir Secret Sharing}} استفاده می‌شود \cite{SecuringVehicularAd}.

ابطال گواهی یکی از دشوارترین ملاحلات است: فهرست ابطال گواهی\LTRfootnote{\lr{Certificate Revocation List (CRL)}} معایبی مانند طولانی بودن، پنجرهٔ آسیب‌پذیری گواهی کوتاه‌عمر و در دسترس نبودن زیرساخت دارد \cite{SecuringVehicularAd}. رایا و هوکسو سه پروتکل \lr{RTPD}\LTRfootnote{\lr{RTPD}} (پیام رمزشدهٔ \lr{CA} به \lr{TPD} که کلیدها را پاک می‌کند و تأیید دریافت می‌فرستد)، \lr{RCCRL}\LTRfootnote{\lr{RCCRL}} (فهرست ابطال فشرده که همسایه‌ها را هم هشدار می‌دهد) و \lr{DRP}\LTRfootnote{\lr{DRP}} (انباشت اتهام توسط همسایه‌ها در حالت خودکار) پیشنهاد می‌کنند و ابطال را «شایداً دشوارترین جنبهٔ امنیت شبکهٔ اقتضایی خودرویی» می‌نامند \cite{SecuringVehicularAd}. مرور فارسیمدان و همکاران می‌نویسند که یک واحد کنارجاده‌ای می‌تواند فهرست ابطال گواهی را در یک منطقهٔ شهری به‌طور کارآمد پخش کند، اما قابلیت اطمینان و ازدحام آن حل‌نشده است \cite{farsimadanReviewSecurityChallenges2025}.

طول عمر گواهی حدود یک روز است و چون مدت رانندگی متغیر است به چند روز کشیده می‌شود؛ با استفادهٔ دو ساعت در روز، حدود ۴۳٬۸۰۰ کلید در سال و حدود ۴٫۲ مگابایت لازم است (با فرض ۱۰۰ بایت برای هر کلید به‌همراه گواهی‌اش) \cite{SecuringVehicularAd}. کاهش این حجم از طریق مشتق کردن کلیدها از کلید اصلی مشترک میان مرجع و خودرو ممکن است، اما باید موازنه‌ای میان اندازهٔ کوچک مجموعهٔ کلید و طول عمر کوتاه گواهی یافت شود \cite{SecuringVehicularAd}.

در مرور ره‌آواتی آگوستینا و همکاران، مرجع قابل اعتماد\LTRfootnote{\lr{Trusted Authority (TA)}} مرجع مرکزی مسئول مدیریت و امنیت ارتباطات در شبکهٔ اقتضایی خودرویی است که با صدور اعتبارنامه‌هایی مانند گواهی دیجیتال و شبه‌نام تضمین می‌کند تنها نهادهای مجاز در شبکه شرکت کنند \cite{rahmawatiagustinaSystematicLiteratureReview2025}. این مرجع نقطهٔ شکست واحد است و عدم در دسترس بودن آن بر همهٔ نهادها اثر می‌گذارد؛ چون هویت همهٔ نهادها را می‌داند، هدف اصلی نقض حریم خصوصی می‌شود و تسخیر آن می‌تواند نقض گستردهٔ حریم خصوصی به همراه داشته باشد \cite{rahmawatiagustinaSystematicLiteratureReview2025}. راه‌حل بررسی‌شده در این مرور، تقسیم مرجع به مرجع ردیابی\LTRfootnote{\lr{Trace Authority (TRA)}} (تولید شبه‌هویت و ردیابی خودروهای مخرب) و مرکز تولید کلید\LTRfootnote{\lr{Key Generation Center (KGC)}} (تولید کلیدهای رمزنگاری و پارامترهای امنیتی) است \cite{rahmawatiagustinaSystematicLiteratureReview2025}.

\subsection{چالش‌های رمزنگاری و امضا}

علی‌رغم این زیرساخت و این انواع امضا، طراحی سازوکارهای رمزنگاری و امضای دیجیتال برای شبکهٔ اقتضایی خودرویی با مجموعه‌ای از چالش‌ها روبه‌روست. نخستین چالش محدودیت منابع محاسباتی واحدهای سوارشونده\LTRfootnote{\lr{On-Board Unit (OBU)}} و کنارجاده‌ای است: توان محاسباتی خودرو کمتر از رایانه، تبلت و گوشی است و عمر طولانی خودرو به این معنا است که سخت‌افزار به‌روز نمی‌شود؛ برخی راه‌حل‌های امنیتی خودرویی نیز به‌دلیل سربار بالا و ظرفیت و عملکرد محاسباتی محدود قابل اجرا نیستند \cite{farsimadanReviewSecurityChallenges2025}. از این رو، الزام دسترس‌پذیری نیازمند استفاده از سازوکارهای رمزنگاری سبک و کم‌هزینه است \cite{farsimadanReviewSecurityChallenges2025}.

چالش دوم سربار تأیید امضا در چگالی بالای ترافیک است: با بزرگ شدن چگالی ترافیک، خودرو نمی‌تواند همهٔ امضاهای پیام‌های همسایگان را به‌موقع تأیید کند که می‌انجامد به از دست رفتن پیام، و سربار ارتباطی نیز در مطالعات پیشین به‌خوبی برطرف نشده است \cite{zhangRAISEEfficientRSUAided2008}. طرح \lr{RAISE}\LTRfootnote{\lr{RAISE}} با کمک واحد کنارجاده‌ای، مسئولیت تأیید اصالت پیام‌های خودروها و اطلاع‌رسانی نتیجه به خودروها را بر عهده دارد و با رویکرد \lr{k}-گمنامی\LTRfootnote{\lr{k-anonymity}} از حریم خصوصی هویت کاربر حفاظت می‌کند؛ بر اساس شبیه‌سازی‌های گستردهٔ گزارش‌شده، این طرح از نظر نسبت از دست رفتن پیام و تأخیر عملکرد بهتری نسبت به طرح‌های پیشین دارد \cite{zhangRAISEEfficientRSUAided2008}.

چالش سوم الزامات زمانی سخت است: محدودیت‌های زمانی به‌خاطر سرعت بالای خودروها از الزامات امنیتی است \cite{SecuringVehicularAd} و ارتباطات با اولویت بالا نباید تأخیری بیش از چند ده میلی‌ثانیه تجربه کنند \cite{farsimadanReviewSecurityChallenges2025}؛ بنابراین عملیات تولید و تأیید امضا نباید به گلوگاه محاسباتی تبدیل شوند. چالش چهارم مدیریت کلید و ابطال گواهی است: در شبکه‌ای با گره‌های متحرک بسیار، صدور، توزیع، به‌روزرسانی و ابطال گواهی پیچیده است و ابطال «شایداً دشوارترین جنبهٔ امنیت» خوانده شده است \cite{SecuringVehicularAd}؛ افزون بر این، در بیشتر سازوکارهای متمرکز واحدهای کنارجاده‌ای به‌عنوان واحدهای کاملاً مورد اعتماد فرض می‌شوند، در حالی که همواره در دسترس نیستند و خودشان آسیب‌پذیرند \cite{farsimadanReviewSecurityChallenges2025}.

چالش پنجم ناتوانی رمزنگاری در برابر حملات داخلی است: راه‌حل‌های رمزنگاری عمدتاً برای دفاع در برابر تهدیدهای خارجی متمرکزند و در برابر اقدامات امنیتی پیشگیرانهٔ داخلی مانند جعل موقعیت و تولید هشدار غلط توسط خودروهای مجاز آسیب‌پذیرند \cite{farsimadanReviewSecurityChallenges2025}. چون پیام‌های ایمنی اطلاعات حساسی ندارند، رمزنگاری برای آن‌ها الزامی نیست و الزام جداگانه‌ای برای بررسی سازگاری و معقول‌بودن داده مطرح شده است \cite{SecuringVehicularAd}. چالش ششم سربار ذاتی رمزنگاری است: سربار رمزنگاری حدود ۱۴۰ بایت است و کارایی امضای دیجیتال جای بهبود دارد \cite{SecuringVehicularAd}. بهینه‌سازی‌های پیشنهادی شامل تأیید پیام فقط در صورت مرتبط بودن، عدم تأیید مجدد کلیدی که پیش‌تر تأیید شده \cite{SecuringVehicularAd}، تأیید دسته‌ای همکارانهٔ امضاها \cite{farsimadanReviewSecurityChallenges2025} و استفاده از رمزنگاری متقارن با افشای تأخیری کلیدها (\lr{TESLA++}\LTRfootnote{\lr{TESLA++}}) و کلیدهای کوتاه‌عمر با تابع هش \cite{farsimadanReviewSecurityChallenges2025} است. استفاده از فناوری‌های دفترکل توزیع‌شده مانند بلاکچین برای تعامل‌پذیری نیز می‌تواند پیچیدگی زمانی اضافه کند \cite{farsimadanReviewSecurityChallenges2025}.

تهدید رایانش کوانتومی نیز در پیشنهادهای آیندهٔ مرور ره‌آواتی آگوستینا مطرح شده است، هرچند فقط در بافت روش تولید شبه‌نام و ناکافی بودن احتمالی توابع هش در بلندمدت \cite{rahmawatiagustinaSystematicLiteratureReview2025}. در مجموع، چالش اصلی در طراحی سازوکارهای رمزنگاری برای این شبکه، تأمین اصالت، یکپارچگی و انکارناپذیری با حداقل سربار محاسباتی و ارتباطی و با رعایت محدودیت‌های زمانی سخت است \cite{SecuringVehicularAd, farsimadanReviewSecurityChallenges2025, zhangRAISEEfficientRSUAided2008}.
